\subsection{Algoritmos de Aproximação com Garantia Formal}
\label{sec:aproximados}

Antes de abordar a busca exata, examinamos os dois algoritmos polinomiais com garantia formal
de fator de aproximação no espaço métrico. Além de constituírem objetos diretos da análise de
complexidade empírica, ambos serão empregados na Seção~\ref{sec:warm_starts} como estratégias
de inicialização para restringir a árvore do \textit{Branch-and-Bound}.

\subsubsection{Algoritmo Twice-Around-The-Tree (\texttt{v0} \texorpdfstring{$\rightarrow$}{->} \texttt{v1})}
\label{sec:twice_around}

\placeholderGrupo{%
  Inserir aqui a descrição metodológica do algoritmo \textit{Twice-Around-The-Tree}
  ($2$-aproximação), o diagnóstico de tempo e memória da implementação inicial
  (\texttt{v0\_naive}), o detalhamento didático das decisões de refatoração adotadas na versão
  otimizada (\texttt{v1\_optimized}, incluindo o compartilhamento do Prim denso $\Theta(N^2)$ e
  a travessia em pré-ordem iterativa em $\Theta(N)$) e as hipóteses teóricas pré-experimentais
  ($H_{TA}$) a serem avaliadas na Seção~\ref{sec:experimentos}.%
}

\subsubsection{Algoritmo de Christofides (\texttt{v0} \texorpdfstring{$\rightarrow$}{->} \texttt{v1})}
\label{sec:christofides}

O algoritmo de Christofides \cite{christofides1976worst} constrói um circuito hamiltoniano de
custo garantidamente limitado a $1{,}5$ vezes o ótimo em qualquer grafo métrico por meio de
quatro estágios encadeados: (1) computar uma Árvore Geradora Mínima (AGM) $T$ sobre o grafo
completo $G$; (2) coletar o conjunto $I \subseteq V$ dos vértices que possuem grau ímpar em
$T$ (cuja cardinalidade $|I|$ é sempre par, pois a soma dos graus de qualquer grafo é o dobro
do número de arestas) e determinar um emparelhamento perfeito de peso mínimo $M$ no subgrafo
completo induzido por $I$ por meio do algoritmo Blossom de Edmonds \cite{edmonds1965paths};
(3) unir as arestas de $T$ e $M$ formando um multigrafo conexo $G' = (V, E(T) \uplus E(M))$ no
qual todo vértice passa a ter grau par, possibilitando extrair um circuito euleriano que
percorre todas as arestas; e (4) percorrer esse circuito eliminando visitas repetidas a
vértices já alcançados (\textit{shortcutting}), operação que não aumenta o custo total graças à
desigualdade triangular.

\paragraph{Diagnóstico da implementação inicial (\texttt{v0\_naive}).}
Uma tradução ingênua direta desses quatro passos como definido no material da disciplina 
sobre bibliotecas genéricas de grafos introduz três
fontes de ineficiência que distorcem tanto o consumo de memória quanto o comportamento
assintótico das subrotinas auxiliares. Primeiro, a versão inicial constrói explicitamente um
grafo completo em memória utilizando dicionários de adjacência dinâmicos contendo todas as
$\binom{N}{2} = N(N-1)/2$ arestas. Em grafos esparsos, listas ou dicionários de adjacência são
vantajosos por armazenarem apenas arestas presentes; em um grafo euclidiano completo, porém,
cada vértice conecta-se a todos os demais, de modo que tabelas de dispersão e objetos de
atributos por aresta consomem $\Theta(N^2)$ memória com uma constante multiplicativa elevada e
alto custo de alocação. Segundo, a extração da AGM recorre ao algoritmo de Kruskal genérico,
que precisa extrair e ordenar todas as $|E| = \Theta(N^2)$ arestas do grafo completo em tempo
$O(|E| \log |E|) = O(N^2 \log N)$. Terceiro, o multigrafo $T \uplus M$ é novamente materializado
em uma estrutura dinâmica de alto nível para obter o circuito euleriano, repetindo alocações
desnecessárias sobre um grafo que possui menos de $1{,}5N$ arestas.

\paragraph{Refatoração gradual para a versão otimizada (\texttt{v1\_optimized}).}
Para eliminar essas sobrecargas e isolar o custo intrínseco de cada etapa teórica, redesenhamos
a arquitetura em quatro decisões coordenadas:

\begin{enumerate}
  \item \textbf{Operação direta sobre matriz de distâncias contígua}: Eliminamos completamente
  a instanciação do grafo completo em estruturas dinâmicas. Todas as consultas de distância
  $c(u, v)$ passam a indexar diretamente a matriz simétrica contígua $N \times N$ pré-computada
  para a instância, reduzindo a localidade de acesso a operações sobre arranjos planos em memória
  e zerando a etapa de construção do grafo.

  \item \textbf{Árvore Geradora Mínima via Prim denso sem fila de prioridades ($\Theta(N^2)$)}:
  Em vez de ordenar todas as $\binom{N}{2}$ arestas como em Kruskal ou manter uma fila de
  prioridades binária (cujo custo $O(|E| \log N)$ também degenera para $O(N^2 \log N)$ em grafos
  completos), implementamos o algoritmo de Prim em sua formulação densa (compartilhada também
  pelo \textit{Twice-Around-The-Tree}). O procedimento mantém quatro vetores unidimensionais de
  tamanho $N$: um indicador booleano de pertencimento à árvore, o vetor \texttt{chave}$[v]$ com
  a menor distância conhecida de cada vértice $v$ até a árvore atual (inicializado com $+\infty$,
  exceto a raiz com $0$), o vetor de predecessores e um vetor de graus na AGM. A cada uma das
  $N$ iterações, o algoritmo varre linearmente o vetor \texttt{chave} em $\Theta(N)$ para
  selecionar o vértice fora da árvore de menor distância $u$; em seguida, marca $u$ como
  integrado à árvore, atribui \texttt{chave}$[u] \leftarrow +\infty$ no próprio vetor e percorre
  a linha $u$ da matriz de distâncias em $\Theta(N)$ atualizando \texttt{chave}$[v]$ para os
  vértices ainda externos sempre que $c(u, v) < \texttt{chave}[v]$.

  Duas escolhas de engenharia merecem destaque nessa estrutura. A marcação de
  \texttt{chave}$[u] = +\infty$ \textit{in-place} garante que vértices já incorporados nunca
  voltem a ser selecionados como mínimos nas iterações futuras, dispensando a criação de
  máscaras booleanas temporárias a cada busca de mínimo. Além disso, sempre que uma aresta
  $(\text{pai}[u], u)$ é confirmada na árvore, incrementam-se imediatamente os graus de ambas as
  extremidades, de modo que o conjunto de vértices ímpares $I$ já sai pronto ao término da AGM
  em uma única varredura linear. Como são executadas exatamente $N$ iterações e cada iteração
  realiza uma busca linear $\Theta(N)$ e uma relaxação linear $\Theta(N)$, o tempo é
  estritamente $\Theta(N^2)$ e o espaço auxiliar estritamente $\Theta(N)$, removendo o fator
  logarítmico de Kruskal e da variante com \textit{heap}.

  \item \textbf{Restrição do algoritmo Blossom ao subgrafo reduzido de vértices ímpares}:
  A biblioteca externa de grafos só é acionada para construir o subgrafo completo induzido pelos
  $|I|$ vértices de grau ímpar e executar o algoritmo Blossom de Edmonds, garantindo
  emparelhamento perfeito mínimo sem instanciar os demais $N - |I|$ vértices pares. Do ponto de
  vista analítico, essa distinção é crucial: a complexidade teórica do algoritmo Blossom é
  $O(|I|^3)$ em tempo e $\Theta(|I|^2)$ em espaço, e não simplesmente uma função homogênea de
  $N$. Para compreender a intuição desse fenômeno, considere duas instâncias com o mesmo número
  total de vértices $N$: se a AGM da primeira tiver formato próximo a um caminho, apenas suas
  duas extremidades terão grau ímpar ($|I| = 2$), tornando o emparelhamento instantâneo; se a
  AGM da segunda apresentar muitas ramificações e folhas (como ocorre em aglomerados ou árvores
  muito estreladas), $|I|$ pode aproximar-se de $N$, elevando drasticamente o esforço do
  Blossom. Logo, o tamanho real $|I|$ do subgrafo ímpar é o parâmetro estrutural que governa o
  custo dessa etapa.

  \item \textbf{Circuito euleriano iterativo via Hierholzer em $\Theta(N)$}:
  Como a AGM contém exatamente $N - 1$ arestas e o emparelhamento $M$ contém $|I|/2 \le N/2$
  arestas, o multigrafo euleriano $G' = (V, E(T) \uplus E(M))$ é sempre esparso, somando
  $|E'| < 1{,}5N = \Theta(N)$ arestas. Representamos $G'$ por listas de adjacência compactas
  nas quais cada aresta inserida recebe um identificador inteiro único de $0$ a $|E'| - 1$,
  acompanhado de um vetor booleano de arestas já consumidas e de um ponteiro de leitura
  monotônico \texttt{ptr}$[u]$ para cada vértice $u$. Atribuir identificadores individuais
  resolve de maneira simples a presença de arestas paralelas (caso uma aresta escolhida em $M$
  coincida com uma aresta já presente em $T$): cada cópia possui seu próprio identificador e é
  marcada independentemente em $\Theta(1)$, sem jamais exigir remoções em meio a listas. O
  algoritmo de Hierholzer avança iterativamente sobre uma pilha explícita (eliminando o risco de
  estouro de pilha de recursão): a partir do vértice no topo $u$, avança \texttt{ptr}$[u]$
  ignorando arestas já marcadas até encontrar uma aresta livre, marca-a e empilha o vizinho;
  quando \texttt{ptr}$[u]$ atinge o final da lista de $u$, desempilha $u$ adicionando-o ao
  circuito euleriano. Como cada entrada de adjacência de $G'$ é inspecionada exatamente uma vez
  pelo ponteiro monotônico ao longo de toda a execução e existem $2|E'| < 3N$ entradas no
  total, tanto o tempo quanto o espaço dessa etapa são estritamente $\Theta(N)$.
\end{enumerate}

\paragraph{Observação de auditoria sobre empates na AGM.}
Durante a verificação formal entre as versões \texttt{v0\_naive} e \texttt{v1\_optimized},
convém notar que em instâncias geométricas com coordenadas em grade regular (como recortes de
\texttt{eil51}), pares distintos de vértices podem apresentar exatamente a mesma distância
euclidiana. Nesses empates exatos, Kruskal e Prim podem selecionar árvores geradoras mínimas
estruturalmente diferentes de mesmo peso total ótimo. Como a mudança de uma única aresta altera
a paridade dos graus de suas extremidades, o conjunto de vértices ímpares $I$ e o ciclo final
obtido após o emparelhamento podem diferir entre as duas versões, mantendo ambas perfeitamente
válidas e sujeitas à mesma garantia formal de $1{,}5$-aproximação.

\paragraph{Hipóteses teóricas pré-experimentais ($H_C$).}
A partir da análise estrutural acima, formulamos três expectativas abstratas a serem validadas
na Seção~\ref{sec:experimentos}: (i) a migração de \texttt{v0\_naive} para
\texttt{v1\_optimized} deve reduzir substancialmente o pico de memória ao eliminar a alocação
quadrática do grafo completo em dicionários; (ii) os tempos medidos de Prim denso e de
Hierholzer devem ajustar-se estreitamente às curvas teóricas $\Theta(N^2)$ e $\Theta(N)$,
respectivamente; e (iii) o emparelhamento Blossom deve concentrar a fração dominante do tempo
total na versão otimizada, apresentando aderência estatística superior quando modelado como
função cúbica de $|I|^3$ (e quadrática em $|I|^2$ para a memória) em comparação ao ajuste
grosseiro em função de $N^3$.
